\chapter{System Analysis}
\label{ch:3}


\begin{figure}[H]
\centering
\includegraphics[width=0.8\textwidth]{02_data_flow.png}
\caption{Data Flow Diagram}
\label{fig:02_data_flow}
\end{figure}


\section{3.1 Functional Requirements}

\subsection{3.1.1 Module 1: Federated Classification (FL Component)}
Requirement FL-C1: Privacy-Preserving Distributed Training

\begin{itemize}
  \item System shall enable 3 hospitals to train a shared classification model without sharing raw patient MRI images
  \item Only model weights (state\_dict) shall be transmitted between clients and central server
  \item Raw patient data shall never leave hospital premises
\end{itemize}

Requirement FL-C2: Multi-Aggregation Strategy Support

\begin{itemize}
  \item System shall support three aggregation algorithms:
\end{itemize}

- FedAvg: Standard weighted averaging baseline

- FedProx: Proximal regularization variant for non-IID data

- QPSO-FL: Quantum-inspired particle swarm optimization (novel method)

\begin{itemize}
  \item Each strategy shall be independently executable with same clients/data for fair comparison
\end{itemize}

Requirement FL-C3: Classification Accuracy

\begin{itemize}
  \item System shall classify brain tumors into 4 classes: Glioma, Meningioma, No Tumor, Pituitary
  \item Minimum global test accuracy target: ≥ 95\%
  \item Per-class performance shall be tracked (precision, recall, F1, ROC-AUC)
\end{itemize}

Requirement FL-C4: Fairness Under Data Heterogeneity

\begin{itemize}
  \item System shall maintain minimum acceptable accuracy for smallest/most data-deprived client (≥ 80\% under moderate label skew)
  \item Fairness metric: Standard deviation of per-client accuracies shall be minimized
  \item System shall explicitly protect minority clients from being sidelined by majority
\end{itemize}

Requirement FL-C5: Convergence and Communication Efficiency

\begin{itemize}
  \item Training shall complete in ≤ 100 communication rounds
  \item Each round shall involve ≤ 5 local training epochs per client
  \item Total training time shall be ≤ 2 hours on standard GPU infrastructure
\end{itemize}


\subsection{3.1.2 Module 2: Segmentation (Local, Non-Federated)}
Requirement SEG-1: Volumetric Segmentation

\begin{itemize}
  \item System shall segment 3D multimodal MRI volumes (4 input channels: T1, T1ce, T2, FLAIR)
  \item Output shall identify 3 tumor sub-regions: Whole Tumor (WT), Tumor Core (TC), Enhancing Tumor (ET)
  \item Segmentation shall be pixel-accurate (voxel-level classification)
\end{itemize}

Requirement SEG-2: Clinical-Grade Accuracy

\begin{itemize}
  \item Mean Dice Score: ≥ 0.75 (clinical-grade threshold)
  \item Tumor Core Dice: ≥ 0.80 (clinically critical region)
  \item Enhancing Tumor Dice: ≥ 0.75
\end{itemize}

Requirement SEG-3: Automated Attention

\begin{itemize}
  \item Model shall use Attention Gates to automatically focus on tumor regions
  \item Healthy tissue suppression shall reduce false positives in non-tumor areas
\end{itemize}

Requirement SEG-4: 3D Volumetric Processing

\begin{itemize}
  \item System shall process full 3D volumes, not 2D slices
  \item Spatial context in 3D shall inform segmentation decisions
\end{itemize}

Requirement SEG-5: Local Processing Only

\begin{itemize}
  \item Segmentation shall occur entirely at individual hospitals
  \item No federated learning, weight sharing, or central aggregation
  \item Each hospital maintains full autonomy over segmentation results
\end{itemize}

KEY DISTINCTION: Module 2 is completely separate from Module 1. Segmentation does NOT use federated learning.


\subsection{3.1.3 Module 3: Progression Forecasting (Local, Non-Federated)}
Requirement PROG-1: Multi-Model Forecasting

\begin{itemize}
  \item System shall fit 4 mathematical models to longitudinal tumor volume data:
\end{itemize}

- Exponential growth

- Gompertz (S-curve with deceleration)

- Logistic (bounded growth)

- Linear (constant rate baseline)

\begin{itemize}
  \item System shall automatically select best-fitting model per patient using R² metric
\end{itemize}

Requirement PROG-2: Hybrid Deep Learning

\begin{itemize}
  \item System shall train LSTM networks to learn residuals (actual - mathematical prediction)
  \item Hybrid model: V\_hybrid = V\_math + LSTM\_correction
  \item Hybrid model shall improve prediction accuracy by ≥ 5\% over baseline
\end{itemize}

Requirement PROG-3: Grade-Stratified Models

\begin{itemize}
  \item Separate models for HGG (High-Grade, aggressive growth) and LGG (Low-Grade, slow growth)
  \item Hyperparameters optimized per grade
\end{itemize}

Requirement PROG-4: Prediction Interval and Confidence

\begin{itemize}
  \item System shall provide point predictions AND uncertainty intervals
  \item 6-month prediction horizon
\end{itemize}

Requirement PROG-5: Local Processing Only

\begin{itemize}
  \item Progression forecasting occurs entirely at individual hospitals
  \item No federated components
  \item Clinical interpretability prioritized over black-box predictions
\end{itemize}

KEY DISTINCTION: Module 3 is completely separate from Module 1. Progression forecasting does NOT use federated learning.


\begin{figure}[H]
\centering
\includegraphics[width=0.85\textwidth]{18_dfd_level0.png}
\caption{DFD Level 0 --- Context Diagram}
\label{fig:18_dfd_level0}
\end{figure}


\section{3.2 Workflow and Clinical Process Flow}

\subsection{3.2.1 Correct Execution Order: Classify → Segment → Forecast}
CRITICAL POINT: Non-Glioma tumors exit at Module 1. This is clinically appropriate:

\begin{itemize}
  \item Meningiomas: Usually extradural, less complex segmentation requirements
  \item Pituitary: Small, standardized treatment often does not require volumetric analysis
  \item No Tumor: No further analysis needed
\end{itemize}


\subsection{3.2.2 Data Flow Diagram (DFD)}
Level 0 (Context Diagram):

See also: Diagram 02\_data\_flow.png and 09\_integration\_architecture.png in /diagrams/rendered/


\begin{figure}[H]
\centering
\includegraphics[width=0.85\textwidth]{09_integration_architecture.png}
\caption{End-to-End Integration Architecture}
\label{fig:09_integration_architecture}
\end{figure}


\section{3.3 Data Requirements}

\subsection{3.3.1 Module 1: Classification Data}
Privacy Consideration: Data never leaves client boundaries. Each hospital maintains full control over its data.


\subsection{3.3.2 Module 2: Segmentation Data}
Local Processing: All segmentation data remains hospital-local. No sharing with federated server.


\subsection{3.3.3 Module 3: Progression Data}
Local Processing: Progression forecasting occurs entirely at individual hospitals.


\section{3.4 Non-Functional Requirements}

\subsection{3.4.1 Performance Requirements}

\subsection{3.4.2 Privacy Requirements}

\subsection{3.4.3 Reliability Requirements}

\subsection{3.4.4 Maintainability Requirements}

\section{3.5 System Constraints}

\subsection{3.5.1 Technical Constraints}
Federated Learning Architecture:

\begin{itemize}
  \item Synchronous updates: All clients must complete local training before aggregation
  \item Central server orchestration: Single server coordinates rounds (simplified vs. peer-to-peer)
  \item No client-to-client communication (simplified privacy model)
\end{itemize}

Scalability Constraints:

\begin{itemize}
  \item Currently tested with 3 clients only
  \item Scaling to 5–10 hospitals untested (computational overhead increases \~O(n))
  \item Communication bandwidth assumes typical hospital internet (≤ model file size ≤ 50MB)
\end{itemize}

Model Capacity Constraints:

\begin{itemize}
  \item ResNet-18: 11.2M parameters (standard, not ultra-lightweight)
  \item Local training assumes ≥ 1 GPU per client or shared GPU time
\end{itemize}


\subsection{3.5.2 Data Constraints}
Classification:

\begin{itemize}
  \item Limited to 2D MRI slices (volumetric processing deferred to Module 2)
  \item Only supports brain-extracted images (preprocessing required)
  \item Fixed input resolution: 224×224
\end{itemize}

Segmentation:

\begin{itemize}
  \item Full 3D volumes required (memory-intensive)
  \item Assumes 4-modality input (some hospitals may only have 1-2 modalities initially)
\end{itemize}

Progression:

\begin{itemize}
  \item Requires ≥ 2 longitudinal timepoints per patient (some patients may have only 1 scan)
  \item Temporal gaps vary (not uniformly spaced)
\end{itemize}


\subsection{3.5.3 Regulatory Constraints}
\begin{itemize}
  \item **HIPAA Compliance:** De-identification required; direct identifiers stripped
  \item **Institutional Review:** Requires IRB approval at each participating hospital
  \item **Data Retention:** Comply with retention policies at each institution
\end{itemize}


\section{3.6 System Dependencies}

\subsection{3.6.1 Software Dependencies}

\subsection{3.6.2 Hardware Dependencies}

\subsection{3.6.3 Deployment Dependencies}
\begin{itemize}
  \item **Kaggle Notebooks:** For initial research (free GPU, simplified setup)
  \item **Local GPU Servers:** For reproducible experiments (NVIDIA GPU + CUDA drivers)
  \item **Streamlit:** Web dashboard for visualization (Python 3.8+)
  \item **Docker:** Optional containerization for multi-hospital deployment
\end{itemize}


\section{3.7 Design Rationale: Why This Architecture?}

\subsection{3.7.1 Why Separate FL Only for Classification?}
Question: Why not federate all three modules?

Answers:

\begin{enumerate}
  \item **Privacy-Utility Trade-off:**
\end{enumerate}

- Classification: Direct output is tumor type (low cardinality, easy to anonymize)

- Segmentation masks: Can leak anatomical patient information (risky to share)

- Progression data: Time-series volumes are sensitive (unique patient signatures)

- FL overhead only justified for high-privacy-value components

\begin{enumerate}
  \item **Computational Efficiency:**
\end{enumerate}

- 3D segmentation requires sliding-window inference (expensive)

- Federated 3D segmentation would require shipping massive 3D volumes or intermediate features (communication bottleneck)

- Local segmentation is more efficient

\begin{enumerate}
  \item **Clinical Workflow:**
\end{enumerate}

- Hospitals already have local segmentation/progression infrastructure

- Classification is the novel collaborative component

- Minimal disruption to existing hospital workflows

\begin{enumerate}
  \item **Regulatory Simplicity:**
\end{enumerate}

- Federal learning adds complexity to compliance

- Hospitals comfort with local-only processing for segments/forecasts

- Federated classification is a pilot, proving value before broader adoption


\subsection{3.7.2 Why QPSO Over FedAvg/FedProx?}
Research Question: What aggregation strategy minimizes clinical inequity?

QPSO Advantages (from literature review + project experiments):

Conclusion: QPSO trades \~1\% global accuracy for massive fairness gains. Clinically, protecting the weakest hospital is more important than marginal global accuracy improvements.


\section{3.8 Chapter Summary}
Figure: Data Flow Diagram - Complete Pipeline


\section{Chapter Summary}
This chapter analyzed system requirements, proposed the three-module architecture, and confirmed technical, operational, and economic feasibility.
